% !TEX root = chapter-standalone.tex

\section{Canonical trace}
\linkvideo{spring2021-par-feedback:feedback} % Feedback
\todotextjira{199}{\alphubel: @JL: Write something about how we already saw string diagrams, and that (co)evaluation maps allow for bending the wires around to change the their direction, and that we can encode dual objects by adding a ``flow'' direction to the wires.
    Then motivate the notion of trace in a \SY{monoidal category} by asking what we get when we bend a wire around and let the outout of a morphism flow back and be it's input -- ``closing the loop''}

%\linkvideo{spring2021-par-feedback:mon-cat:string-diag} % String diagrams

\begin{ctdefinition}[Trace of an endomorphism]
    \label{def:trace_endo}
    Let~$\tup{\CatC, \mtimescat, \idmoncat, \associator, \leftunitor, \rightunitor, \braiding}$ be a \SY{symmetric monoidal category}.
    Let~$\Obja \setin \ObC$ be dualizable and let
    \begin{equation}
        \mora \colon \Obja \mto \Obja.
    \end{equation}
    The \emph{trace} of~$\mora$ is the morphism
    \begin{equation}
        \Tr(\mora) \colon \idmoncat \mto \idmoncat
    \end{equation}
    defined by
    \equationsag{endo_trace}{eq:endo_trace}
    %\begin{equation}
    %\label{eq:endomorphism-gen-trace}
    %\idmoncat \overset{\coev_\Obja}{\mto} \Obja \mtimescatat \Obja^\vee  \overset{\mora \mtimescatmor \catid_{\Obja^\vee}}{\mto} \Obja \mtimescatob \Obja^\vee  \overset{\braiding}{\mto}  \Obja^\vee \mtimescatob \Obja \overset{\ev_\Obja}{\mto} \idmoncat.
    %\end{equation}
\end{ctdefinition}

\begin{gradedexercise}[\exname{LinearAlgebraTrace}]
    \label{ex:LinearAlgebraTrace}
    Let~\CatC be the category of finite-dimensional real \SY{vector spaces} and~\reals-linear maps.
    \todotextjira{516}{\alphubel: @JL: Didn't we call this already something different?}
    We have seen that this category is symmetric monoidal when equipped with the usual tensor product as monoidal product.
    Furthermore, in \cref{sec:dual-objects} we saw that every object in this category is dualizable.

    Fix a finite-dimensional real \SY{vector space}~$\styleobj{V}$, and let~$\makeset{e_1,\ldots,e_n }$ be a basis of it.
    We saw that a choice of dual object for~$\styleobj{V}$ is given by~$\styleobj{V}^* = \Hom(\styleobj{V}, \reals)$, together with the evaluation map
    \begin{equation}
        \defmapperiod{
            \ev_\styleobj{V}
        }{
            \styleobj{V}^* \otimes \styleobj{V}
        }{
            \mto
        }{
            \reals
        }{
            \tup{l, v}
        }{
            l(v)
        }
    \end{equation}
    and the co-evaluation map
    \begin{equation}
        \defmapperiod{
            \coev_\styleobj{V}
        }{
            \reals
        }{
            \mto
        }{
            \styleobj{V} \otimes \styleobj{V}^*
        }{
            \lambda
        }{
            \lambda \sum_{i=1}^n e_i \otimes e_i^*
        }
    \end{equation}
    where $\makeset{e_1^*,\ldots,e_n^* }$ is the basis dual to the one we chose for~$\styleobj{V}$.

    Let~$\mora\colon \styleobj{V} \mto \styleobj{V}$ be a linear \SY{endomorphism} -- that is,~$\mora \setin \Hom_{\CatC}(\styleobj{V}, \styleobj{V})$.
    Compute the trace~$\Tr(\mora) \setin \Hom_{\CatC}(\reals, \reals)$ of~$\mora$ according to \cref{def:trace_endo}, and explain why it is the linear map ``multiplication by the trace of~$\mora$'', where ``trace'' in this latter phrase is the usual notion that we know from linear algebra.
\end{gradedexercise}
\solutionof{LinearAlgebraTrace}

\begin{gradedexercise}[\exname{DPSnakeTrace}]
    \label{ex:DPSnakeTrace}

    In this exercise we work with the category $\DP$ of posets and design problems, equipped with the symmetric monoidal structure where the monoidal product is the cartesian product of posets.
    The braiding
    $$\braiding_{\posgenA,\posgenB}\colon \F{\posgenA} \Ptimes \F{\posgenB} \profto \R{\posgenB} \Ctimes \R{\posgenA}$$
    is defined by
    \begin{equation}
        \braiding_{\posgenA,\posgenB}(\tup{\F{\posgenAel_1},\F{\posgenBel_1}}\Fop, \tup{\R{\posgenBel_2},\R{\posgenAel_2}})\definedas \pars{\F{\posgenAel_1}\posleqof\posA \R{\posgenAel_2}}\booland \pars{\F{\posgenBel_1}\posleqof\posB \R{\posgenBel_2}}.
    \end{equation}
    In the following you are free to use the identification
    \begin{equation}
        (\posgenA \Ptimes \posgenB)\op = \posgenA\op \Ptimes \posgenB\op
    \end{equation}
    for any posets $\posgenA$, $\posgenB$.
    Also, recall that $(\posgenA\op)\op = \posgenA$.

    Let us define the following duality data:
    \begin{itemize}
        \item $\posgenA\dual \definedas \posgenA\op$
        \item $\ev_\posgenA \colon \F{(\posgenA\op \Ptimes \posgenA)}\op \Ptimes \R{\makeset{\bullet}} \mto \Bool, \quad \tup{\tup{\ela^*, \elb}^*, \bullet} \mapsto \elb \posleqof\posgenA \ela$
        \item $\coev_\posgenA \colon \F{\makeset{\bullet}}\op \Ptimes \R{(\posgenA \Ptimes \posgenA\op)} \mto \Bool, \quad \tup{\bullet, \tup{\ela, \elb^*}} \mapsto \elb \posleqof\posgenA \ela$
    \end{itemize}

    Your tasks:

    \begin{enumerate}
        \item
              Guess the definitions of the associator $\associator$ and the unitors $\leftunitor$, $\rightunitor$ for the monoidal category $\DP$, check that each has the correct type, and justify why each of them does indeed define a morphism in the category $\DP$.
        \item
              Guess the definitions of $\leftunitor^{-1}$ and $\rightunitor^{-1}$ and check for one of them that it does indeed define the inverse morphism.
        \item
              Check that $\ev_\posgenA$ and $\coev_\posgenA$, as defined above, are morphisms in $\DP$.
        \item
              For an arbitrary poset $\posgenA$ and the duality data given above, prove that this snake equation
              \begin{equation}
                  \leftunitor_{\posgenA}^{-1} \mthen (\coev_\posgenA \mtimescatmor \catidat \posgenA) \mthen \associator_{\posgenA, \posgenA\op, \posgenA} \mthen(\catidat\posgenA \mtimescatmor \ev_\posgenA) \mthen \rightunitor_{\posgenA} = \catidat\posgenA
              \end{equation}
              holds.
    \end{enumerate}
\end{gradedexercise}
\solutionof{DPSnakeTrace}

\begin{widepar}
    \begin{ctdefinition}[Trace of a generalized endomorphism]
        \label{def:trace_gen_endo}
        Let~$\tup{\CatC, \mtimescat, \idmoncat, \associator, \leftunitor, \rightunitor, \braiding}$ be a \SY{symmetric monoidal category}.
        Let~$\Obja \setin \ObC$ be dualizable and let
        \begin{equation}
            \mora \colon (\Objb \mtimescatob \Obja) \mto( \Objc \mtimescatob \Obja).
        \end{equation}
        The \emph{trace over}~$\Obja$ of~$\mora$ is the morphism
        \begin{equation}
            \Tr_{\Objb, \Objc}^\Obja (\mora) \colon \Objb \mto \Objc
        \end{equation}
        defined by
        \equationsag{endo_trace_gen}{eq:endo_trace_gen}
    \end{ctdefinition}
\end{widepar}

\subsection{Canonical trace for linear relations}

We now show that the symmetric monoidal category of finite-dimensional real vector spaces and linear relations, with direct sum as monoidal product, is \SY{compact closed}, and therefore admits a canonical trace.

We consider the category~$\stylecat{LinRel}_\reals$ (\cref{def:cat-of-linear-relations}) whose objects are finite-dimensional real vector spaces and whose morphisms are linear relations (i.e., a morphism from~$\styleobj{V}$ to~$\styleobj{W}$ is a linear subspace~$\relA \subseteq \styleobj{V} \oplus \styleobj{W}$).
This category is symmetric monoidal with the direct sum~$\oplus$ as monoidal product and the zero vector space~$\makeset{0}$ as monoidal unit.

\subsubsection{Duality in~$\stylecat{LinRel}_\reals$}

We claim that every object in~$\stylecat{LinRel}_\reals$ is \SY{right dualizable}, and therefore the category is \SY{compact closed}.
In fact, the dual of any vector space is the space itself: for any finite-dimensional real vector space~$\styleobj{V}$, we set
\begin{equation}
    \styleobj{V}\dual = \styleobj{V}.
\end{equation}
The evaluation and coevaluation are defined as follows.
The evaluation
\begin{equation}
    \ev_\styleobj{V} \colon \styleobj{V} \oplus \styleobj{V} \mto \makeset{0}
\end{equation}
is the linear relation
\begin{equation}
    \ev_\styleobj{V} = \makeset{ \tup{\tup{v, w}, 0} \setin (\styleobj{V} \oplus \styleobj{V}) \oplus \makeset{0} \mid v + w = 0 }.
\end{equation}
In other words,~$\tup{v,w}$ is related to~$0$ if and only if~$w = -v$.

The coevaluation
\begin{equation}
    \coev_\styleobj{V} \colon \makeset{0} \mto \styleobj{V} \oplus \styleobj{V}
\end{equation}
is the linear relation
\begin{equation}
    \coev_\styleobj{V} = \makeset{ \tup{0, \tup{v, w}} \setin \makeset{0} \oplus (\styleobj{V} \oplus \styleobj{V}) \mid v + w = 0 }.
\end{equation}

\begin{remark}
    One can verify that~$\ev_\styleobj{V}$ and~$\coev_\styleobj{V}$ are indeed linear subspaces (of the appropriate direct sums), and that the snake equations (\cref{eq:snake-eq-1,eq:snack-eq-2}) are satisfied.
    The key point is that the condition~$v + w = 0$ defines a linear subspace, and the composition of these relations with the associators and unitors (which in~$\stylecat{LinRel}_\reals$ with~$\oplus$ are the obvious identifications) yields the identity relation.
\end{remark}

\subsubsection{The canonical trace formula for linear relations}

Since~$\stylecat{LinRel}_\reals$ is \SY{compact closed}, the canonical trace (\cref{def:trace_gen_endo}) is defined.
Given a linear relation
\begin{equation}
    \relA \subseteq (\styleobj{V} \oplus \styleobj{W}) \oplus (\styleobj{V'} \oplus \styleobj{W})
\end{equation}
viewed as a morphism~$\relA \colon \styleobj{V} \oplus \styleobj{W} \mto \styleobj{V'} \oplus \styleobj{W}$ in~$\stylecat{LinRel}_\reals$, the trace over~$\styleobj{W}$ is the linear relation
\begin{equation}
    \Tr_{\styleobj{V}, \styleobj{V'}}^{\styleobj{W}}(\relA) \colon \styleobj{V} \mto \styleobj{V'}
\end{equation}
given by
\begin{equation}
    \label{eq:canonical-trace-linrel}
    \Tr_{\styleobj{V}, \styleobj{V'}}^{\styleobj{W}}(\relA) = \makeset{ \tup{v, v'} \setin \styleobj{V} \oplus \styleobj{V'} \mid \exists\, w \setin \styleobj{W} \colon \tup{\tup{v, w}, \tup{v', w}} \setin \relA }.
\end{equation}

In other words, the trace identifies the ``input copy'' and the ``output copy'' of~$\styleobj{W}$: one asks for the existence of a vector~$w \setin \styleobj{W}$ that serves simultaneously as the~$\styleobj{W}$-component of both the input and the output.

\begin{remark}
    Comparing \cref{eq:canonical-trace-linrel} with the trace formula for~$\Rel$ with cartesian product (\cref{sec:dual-objects}), we see the same pattern: the trace is obtained by ``identifying the feedback wire'', i.e., requiring that the component in~$\styleobj{W}$ on the input side equals (in~$\Rel$) or sums to zero with (in the evaluation for~$\stylecat{LinRel}_\reals$) the component on the output side.
    In the present case, the condition~$v + w = 0$ in the (co)evaluation means that the feedback wire carries~$w$ on one end and~$-w$ on the other; after composing, this yields the condition that the same~$w$ appears in both slots of~$\relA$.
\end{remark}

